% IEEE Access — figure includes
% Generated automatically

\begin{figure}[htbp]
  \centering
  \includegraphics[width=\columnwidth]{figures/fig01_pipeline_overview.pdf}
  \caption{End-to-end TDE pipeline from ESRF CT acquisition to labeled synthetic deformation dataset. Stages correspond to the three processing modules: data preparation, 3-D projection, and synthetic deformation generation.}
  \label{fig:pipeline_overview}
\end{figure}

\begin{figure}[htbp]
  \centering
  \includegraphics[width=\columnwidth]{figures/fig02_module_dependency.pdf}
  \caption{Software module dependency graph of the TDE pipeline. Nodes represent Python modules; directed edges denote import or data-flow dependencies.}
  \label{fig:module_dependency}
\end{figure}

\begin{figure}[htbp]
  \centering
  \includegraphics[width=\columnwidth]{figures/fig03_data_flow.pdf}
  \caption{Data flow diagram of the TDE pipeline showing transformations applied at each processing step and the intermediate data formats.}
  \label{fig:data_flow}
\end{figure}

\begin{figure}[htbp]
  \centering
  \includegraphics[width=\columnwidth]{figures/fig04_specimen_distribution.pdf}
  \caption{Number of 2-D projection images generated per fossil specimen. Each specimen contributes exactly 30 projections (6 orthographic + 24 rotational), yielding a balanced dataset.}
  \label{fig:specimen_distribution}
\end{figure}

\begin{figure}[htbp]
  \centering
  \includegraphics[width=\columnwidth]{figures/fig05_dataset_composition.pdf}
  \caption{Dataset composition showing equal contribution of all eight archosaur specimens to the TDE synthetic deformation dataset. Each specimen contributes 120 synthetic images (12.5%).}
  \label{fig:dataset_composition}
\end{figure}

\begin{figure}[htbp]
  \centering
  \includegraphics[width=\columnwidth]{figures/fig06_dataset_growth.pdf}
  \caption{Dataset scale at each stage of the TDE pipeline, demonstrating the multiplicative effect of multi-view projection (×30) and synthetic deformation (×4) on the base specimen count.}
  \label{fig:dataset_growth}
\end{figure}

\begin{figure}[htbp]
  \centering
  \includegraphics[width=\columnwidth]{figures/fig07_projection_types.pdf}
  \caption{Breakdown of the 30 projection views generated per specimen: 6 orthographic views (3 axes × 2 directions) and 24 rotational views (0° to 345° at 15° increments) of the sagittal MIP.}
  \label{fig:projection_types}
\end{figure}

\begin{figure}[htbp]
  \centering
  \includegraphics[width=\columnwidth]{figures/fig08_deformation_distribution.pdf}
  \caption{Distribution of synthetic images across four taphonomic deformation categories. All categories are balanced with 240 images each, producing a class-balanced training dataset.}
  \label{fig:deformation_distribution}
\end{figure}

\begin{figure}[htbp]
  \centering
  \includegraphics[width=\columnwidth]{figures/fig09_projection_pipeline.pdf}
  \caption{Streaming Maximum Intensity Projection (MIP) architecture. JP2 slices are decoded one at a time; three orthogonal MIPs are accumulated in a single pass, reducing peak memory from O(D \cdot H^2) to O(H^2).}
  \label{fig:projection_pipeline}
\end{figure}

\begin{figure}[htbp]
  \centering
  \includegraphics[width=\columnwidth]{figures/fig10_sample_projections.pdf}
  \caption{Representative 2-D Maximum Intensity Projection images from four archosaur specimens. Top row: orthographic coronal views. Bottom row: rotational views (45°, 90°, 135°, 180°). All images are 224×224 grayscale PNG.}
  \label{fig:sample_projections}
\end{figure}

\begin{figure}[htbp]
  \centering
  \includegraphics[width=\columnwidth]{figures/fig11_all_specimens_grid.pdf}
  \caption{Complete gallery of orthographic coronal (ortho\_00) Maximum Intensity Projection images for all eight archosaur specimens in the TDE dataset. Images are 224×224 grayscale PNG.}
  \label{fig:all_specimens_grid}
\end{figure}

\begin{figure}[htbp]
  \centering
  \includegraphics[width=\columnwidth]{figures/fig12_voxel_sizes.pdf}
  \caption{CT voxel resolution (µm) for each archosaur specimen. Voxel sizes span two orders of magnitude (1.28–107.16 µm), reflecting the diverse scan resolutions available in the ESRF Paleo CT dataset.}
  \label{fig:voxel_sizes}
\end{figure}

\begin{figure}[htbp]
  \centering
  \includegraphics[width=\columnwidth]{figures/fig13_transformation_matrices.pdf}
  \caption{Transformation matrix representations for the four taphonomic deformations. (a) Compression: $c \in [0.3, 0.9]$. (b) Shearing: $k \in [-0.5, 0.5]$. (c) Stretching: $s_x, s_y \in [1.0, 1.5]$. (d) Dissolution: intensity subtraction via ellipse mask.}
  \label{fig:transformation_matrices}
\end{figure}

\begin{figure}[htbp]
  \centering
  \includegraphics[width=\columnwidth]{figures/fig14_deformation_geometry.pdf}
  \caption{Geometric illustration of the four taphonomic deformation types applied to a unit square. The dashed square shows the original boundary; the solid polygon shows the deformed result.}
  \label{fig:deformation_geometry}
\end{figure}

\begin{figure}[htbp]
  \centering
  \includegraphics[width=\columnwidth]{figures/fig15_deformation_gallery.pdf}
  \caption{Deformation gallery for the araripesaurus specimen (ortho\_00 view). Each column shows: original projection, deformed image, and binary deformation mask. Rows correspond to compression, shearing, stretching, and dissolution respectively.}
  \label{fig:deformation_gallery}
\end{figure}

\begin{figure}[htbp]
  \centering
  \includegraphics[width=\columnwidth]{figures/fig16_parameter_distributions.pdf}
  \caption{Parameter distributions for all four deformation types across the full TDE dataset (n=240 per type, seed=42). Parameters are uniformly sampled: $c \sim \mathcal{U}(0.3,0.9)$, $k \sim \mathcal{U}(-0.5,0.5)$, $s_x, s_y \sim \mathcal{U}(1.0,1.5)$, $\lambda \sim \mathcal{U}(0.3,0.8)$.}
  \label{fig:parameter_distributions}
\end{figure}

\begin{figure}[htbp]
  \centering
  \includegraphics[width=\columnwidth]{figures/fig17_stretching_scatter.pdf}
  \caption{Joint distribution of stretching parameters $s_x$ and $s_y$ (n=240). Both are independently sampled from $\mathcal{U}(1.0,1.5)$, confirming near-zero correlation and independent axis scaling.}
  \label{fig:stretching_scatter}
\end{figure}

\begin{figure}[htbp]
  \centering
  \includegraphics[width=\columnwidth]{figures/fig18_compression_progression.pdf}
  \caption{Effect of increasing compression factor $c$ on specimen appearance. A lower value of $c$ produces stronger vertical compression, simulating sedimentary overburden deformation.}
  \label{fig:compression_progression}
\end{figure}

\begin{figure}[htbp]
  \centering
  \includegraphics[width=\columnwidth]{figures/fig19_dissolution_masks.pdf}
  \caption{Example binary deformation masks for the dissolution deformation type. Each mask encodes the elliptical region where intensity was attenuated. Masks are provided for Grad-CAM explainability validation.}
  \label{fig:dissolution_masks}
\end{figure}

\begin{figure}[htbp]
  \centering
  \includegraphics[width=\columnwidth]{figures/fig20_mask_coverage_stats.pdf}
  \caption{Distribution of deformation mask coverage (\% of image pixels changed) per deformation type. Dissolution masks reflect the random ellipse area; compression, shearing, and stretching masks reflect the affine-transformed region.}
  \label{fig:mask_coverage}
\end{figure}

\begin{figure}[htbp]
  \centering
  \includegraphics[width=\columnwidth]{figures/fig21_intensity_distributions.pdf}
  \caption{Pixel intensity distributions (0–255) for original projection images and each synthetic deformation type, sampled from araripesaurus\_\_BSPG-1982-I-90. Dissolution is the only deformation that shifts the mean intensity downward.}
  \label{fig:intensity_distributions}
\end{figure}

\begin{figure}[htbp]
  \centering
  \includegraphics[width=\columnwidth]{figures/fig22_folder_structure.pdf}
  \caption{Canonical folder hierarchy of the TDE dataset. Each specimen directory contains an ESRF-standard \texttt{org\_slices/} directory with JP2 CT slices and a \texttt{mesh/} directory for 3-D surface meshes.}
  \label{fig:folder_structure}
\end{figure}

\begin{figure}[htbp]
  \centering
  \includegraphics[width=\columnwidth]{figures/fig23_function_call_graph.pdf}
  \caption{Key function call graph for the TDE projection and deformation pipelines. Edge direction indicates caller → callee. Nodes are colored by processing module.}
  \label{fig:function_call_graph}
\end{figure}

\begin{figure}[htbp]
  \centering
  \includegraphics[width=\textwidth]{figures/fig24_4deformation_panel.pdf}
  \caption{Publication-quality comparison panel of the four taphonomic deformations applied to the Archaeopteryx\_london specimen (BMNH-37001, ortho\_00). Each row shows original projection, synthetic deformed image, and pixel-wise binary mask.}
  \label{fig:4deformation_panel}
\end{figure}

\begin{figure}[htbp]
  \centering
  \includegraphics[width=\columnwidth]{figures/fig25_rotational_strip.pdf}
  \caption{Twelve representative rotational MIP views of the halszkaraptor specimen at 30° increments. Views are generated by rotating the sagittal MIP (axis-1 projection) in 2-D after accumulation.}
  \label{fig:rotational_strip}
\end{figure}

\begin{figure}[htbp]
  \centering
  \includegraphics[width=\columnwidth]{figures/fig26_mip_algorithm.pdf}
  \caption{Conceptual illustration of the streaming Maximum Intensity Projection (MIP) algorithm. JP2 slices are processed one at a time; a running \texttt{np.maximum} accumulates the global MIP in O(H·W) memory regardless of volume depth D.}
  \label{fig:mip_algorithm}
\end{figure}

\begin{figure}[htbp]
  \centering
  \includegraphics[width=\columnwidth]{figures/fig27_percentile_normalization.pdf}
  \caption{Percentile normalization (1--99\%) applied to a raw MIP accumulator output from the araripesaurus specimen. Left: raw uint16 histogram; Right: normalized uint8 output after clipping and rescaling.}
  \label{fig:percentile_normalization}
\end{figure}

\begin{figure}[htbp]
  \centering
  \includegraphics[width=\columnwidth]{figures/fig28_taxonomy_voxel.pdf}
  \caption{Archosaur specimen taxonomy and CT voxel resolution. Marker size encodes voxel resolution (µm). The dataset spans birds (Aves), non-avian dinosaurs (Dromaeosauridae), and flying reptiles (Pterosauria).}
  \label{fig:taxonomy_voxel}
\end{figure}

\begin{figure}[htbp]
  \centering
  \includegraphics[width=\textwidth]{figures/fig29_dataset_summary.pdf}
  \caption{TDE dataset summary statistics. The dataset contains 240 projection images and 960 synthetic images with corresponding binary masks and per-image labels across four balanced deformation classes.}
  \label{fig:dataset_summary}
\end{figure}

\begin{figure}[htbp]
  \centering
  \includegraphics[width=\columnwidth]{figures/fig30_csv_schema.pdf}
  \caption{Schema of the \texttt{synthetic\_labels.csv} output file. Each row describes one synthetic image with its specimen, deformation type, sampled parameters, output filenames, and corresponding binary mask filename.}
  \label{fig:csv_schema}
\end{figure}
